\section{The Staircase to \texorpdfstring{$\{3,4,3,4\}$}{3,4,3,4}}

No one found $\{3,4,3,4\}$ by seeing it. It was reached by slowly recognizing deeper structural patterns in symmetry, regularity, reflection, curvature, and dimension. The path from a checkerboard to $\{3,4,3,4\}$ is essentially the history of geometry unfolding, and the lesson it teaches is the one the whole theory rests on. This geometry is forced, not chosen.

\subsection{The conceptual staircase}

Each step generalized exactly one thing. First adjacency, then symmetry, then recursion, then dimension, then curvature, and finally ideal infinity.

\begin{enumerate}
  \item A checkerboard.
  \item Regular tilings, of which there are only three flat ones, $\{3,6\}$, $\{4,4\}$, and $\{6,3\}$.
  \item The Platonic solids, where local angle constraints fix the global structure.
  \item Recursive regularity, captured by the Schlafli symbol, so that a cube is $\{4,3\}$.
  \item Higher dimensions, where the polytopes of 4D, the polychora, appear.
  \item The $24$-cell, an exceptional 4D-only shape with no 3D analogue.
  \item Reflection groups, where geometry is generated from mirrors rather than drawn.
  \item Coxeter diagrams, where a tiny integer diagram encodes a whole geometry.
  \item Curvature freedom, where angle sums force a spherical or hyperbolic world.
  \item Hyperbolic space, where the parallel postulate becomes optional and room explodes.
  \item Ideal structures, the vertices at infinity, horospheres, and cusps.
  \item Hyperbolic 4-space honeycombs.
  \item The cubic-honeycomb boundary, since the vertex figure of $\{3,4,3,4\}$ is $\{4,3,4\}$, flat 3D space.
  \item The honeycomb $\{3,4,3,4\}$.
\end{enumerate}

The decisive milestones are easy to name. Schlafli's higher-dimensional polytopes and the $24$-cell came in the 1850s. The discovery of hyperbolic geometry by Lobachevsky, Bolyai, and Gauss came in the 1820s through 1840s. The reflection-group viewpoint and Coxeter's grand unification arrived in the 1930s, showing that regular polytopes, reflection groups, hyperbolic geometry, and higher dimensions are all one recursive symmetry language \cite{coxeter1973}. The realization that the ideal corner of $\{3,4,3,4\}$ is the ordinary cubic lattice, Euclidean 3-space appearing on the boundary of hyperbolic 4-space, is the structural payoff the whole staircase climbs toward.

\subsection{The hidden idea}

The deepest insight in the whole climb is that geometry can be generated algebraically. Not by drawing shapes by hand but by symmetry generators, reflection relations, and combinatorial rules, the integers of the Schlafli symbol. That changed geometry permanently, and it is exactly how the substrate is built. The crystal is the reflection orbit of one chamber, a finite rule unfolding forever.

\begin{result}[Geometry from a finite rule]
The whole of $\{3,4,3,4\}$, infinite and four-dimensional, is the orbit of a single fundamental chamber under five reflections. A finite local rule generates the entire large-scale structure, just as the model's dynamics generate physics from a single local update.
\end{result}

\subsection{Why it feels discovered, not invented}

These structures are rigid, highly constrained, and mathematically forced. Nobody arbitrarily invented $\{3,4,3,4\}$. Once the rules of regularity and curvature are accepted, it almost reveals itself, as one of the maximally symmetric ideal extensions of $24$-cell geometry into negatively curved infinite space. This discovered-not-invented character is the theory's central claim about its own substrate. The selection made later in the paper does not pick $\{3,4,3,4\}$ from a menu. The requirements, a crystallographic spinor coin, curvature, and a 3D cusp, force it.

The staircase is the same logic the theory runs on, local repeating rules generating large-scale order, geometry from reflection, and infinite hyperbolic order emerging naturally from local laws. The substrate is not designed shape by shape. It is the inevitable unfolding of a small set of constraints, with $\{3,4,3,4\}$ standing as one grand culmination of that chain.
